% © 2025 Ing. David Jaroš — CC BY-NC-ND 4.0
%
% This work is licensed under a Creative Commons Attribution-NonCommercial-NoDerivatives
% 4.0 International License (CC BY-NC-ND 4.0).
%
% License History: Earlier drafts (up to v0.3) were released under CC BY 4.0.
% From v0.4 onward, all material is released under CC BY-NC-ND 4.0 to protect
% the integrity of the theoretical work during ongoing academic development.

\documentclass[12pt]{article}
\usepackage{amsmath,amssymb,amsthm}
\usepackage{geometry}
\usepackage{hyperref}
\usepackage{tcolorbox}
\usepackage{xcolor}
\geometry{margin=1in}

% Conjecture environment
\newtheorem{conjecture}{Conjecture}
\newtheorem{definition}{Definition}
\newtheorem{remark}{Remark}

\title{ST-1: Causal Structure of Complex Time $\tau = t + i\psi$ in\\
Unified Biquaternion Theory\\[0.5em]
\large\textsc{Speculative Extension}}
\author{David Jaroš}
\date{2026-03-01}

\begin{document}
\maketitle

\begin{tcolorbox}[colback=yellow!10, title=Status]
\textbf{SPECULATIVE EXTENSION} — This content is not part of
the canonical UBT framework. It is an exploratory investigation
and has not been validated against the canonical field equations.
\end{tcolorbox}

\vspace{1em}
\noindent\textit{Output of Sub-task ST-1 from the CTC/Causal-Structure investigation track.
See also: \texttt{rotation\_complex\_time.tex}, \texttt{ctc\_conditions.tex},
\texttt{chronology\_protection\_ubt.tex}, \texttt{ctc\_methodology.tex}.}

%----------------------------------------------------------------------
\section*{Abstract}
%----------------------------------------------------------------------
We define the causal structure of Unified Biquaternion Theory (UBT) with
complex time $\tau = t + i\psi$.  Three candidate interpretations of the
imaginary component $\psi$ are examined: (i) second independent time
direction, (ii) internal phase coordinate, and (iii) auxiliary degree of
freedom.  A causal partial order on the complex time plane is proposed and
the physical consequences of each interpretation are explored.
All conclusions are \textbf{speculative}.

%----------------------------------------------------------------------
\section{Motivation and Setup}
%----------------------------------------------------------------------
UBT introduces complex time
\begin{equation}
  \tau \;=\; t + i\psi , \qquad t,\,\psi \in \mathbb{R},
  \label{eq:complex_time}
\end{equation}
where $t$ is the familiar Lorentzian (real) time coordinate and $\psi$ is
a phase coordinate that appears in the imaginary sector of the biquaternionic
metric
\begin{equation}
  \mathcal{G}_{\mu\nu} \;=\; g_{\mu\nu} + i\,h_{\mu\nu}
                                + j\,a_{\mu\nu} + k\,b_{\mu\nu}, \qquad
  g_{\mu\nu} := \operatorname{Re}(\mathcal{G}_{\mu\nu}).
  \label{eq:biquat_metric}
\end{equation}
Standard GR operates entirely on the real sector $g_{\mu\nu}$.  Any causal
structure in UBT must reduce to standard Lorentzian causality when
$\psi \to 0$.

%----------------------------------------------------------------------
\section{Candidate Interpretations of $\psi$}
%----------------------------------------------------------------------
\subsection{Interpretation A: Second Independent Time Direction}

If $\psi$ is treated as a genuinely independent time coordinate alongside
$t$, the spacetime acquires an effective $(1+1,3)$ signature — two time-like
and three spatial dimensions.  This is the most radical reading.

\paragraph{Physical consequences.}
\begin{itemize}
  \item The causal cone splits into two orthogonal temporal axes: ordinary
    Lorentzian light-cones (driven by $t$) and a new family of
    $\psi$-directed cones.
  \item Closed paths in the $(t,\psi)$ plane can wrap around the $\psi$
    axis and return to their starting point even if $t$ advances
    monotonically — a potential CTC mechanism.
  \item Standard chronology-protection arguments (Hawking \cite{Hawking1992})
    would need to be re-derived in this extended signature.
  \item The unitarity of the $\psi$-direction evolution is not guaranteed
    by the real-sector field equations and must be imposed as a separate
    constraint.
\end{itemize}

\subsection{Interpretation B: Internal Phase Coordinate (Preferred)}

Under this interpretation $\psi$ parametrises an internal phase degree of
freedom, analogous to a compact extra dimension with period $2\pi R_\psi$.
Real-time $t$ remains the sole causal variable.

\paragraph{Physical consequences.}
\begin{itemize}
  \item Causal order is determined exclusively by the real metric $g_{\mu\nu}$
    (as in standard GR), so no new causal pathologies appear at leading order.
  \item The imaginary sector $h_{\mu\nu},a_{\mu\nu},b_{\mu\nu}$ contributes
    to the stress-energy through the biquaternionic field equation
    $\mathcal{E}_{\mu\nu} = \kappa\,\mathcal{T}_{\mu\nu}$ and can
    deform light cones only through back-reaction on $g_{\mu\nu}$.
  \item A Wick rotation $t \to -i\psi$ is a formal mathematical tool and
    does not carry physical time-reversal significance.
  \item This interpretation is most consistent with the canonical UBT
    programme.
\end{itemize}

\subsection{Interpretation C: Auxiliary Degree of Freedom}

$\psi$ could be purely auxiliary — a Lagrange multiplier or a gauge
degree of freedom eliminated by a constraint.  In this case $\psi$ carries
no dynamical or causal content at all.

\paragraph{Physical consequences.}
\begin{itemize}
  \item All causal structure collapses to standard GR causality.
  \item No new CTC mechanisms arise from $\psi$ alone.
  \item The imaginary-sector equations become constraint equations rather
    than evolution equations.
\end{itemize}

%----------------------------------------------------------------------
\section{Proposed Causal Order on the Complex Time Plane}
%----------------------------------------------------------------------
\begin{definition}[Complex-time causal order]
Let $\tau_1 = t_1 + i\psi_1$ and $\tau_2 = t_2 + i\psi_2$.  Under
Interpretation B, the \emph{causal order} on $\mathbb{C}_\tau$ is
\begin{equation}
  \tau_1 \preceq \tau_2 \quad\Longleftrightarrow\quad
  t_1 \leq t_2 \quad\text{and}\quad |\psi_2 - \psi_1| \leq \Lambda_\psi(t_2-t_1),
\end{equation}
where $\Lambda_\psi > 0$ is a characteristic rate that bounds how quickly
the phase coordinate can vary relative to real time.  The relation
$\preceq$ is reflexive and transitive.
\end{definition}

\begin{remark}
If $\Lambda_\psi \to \infty$, the $\psi$ coordinate is unconstrained and
the order reduces to $t_1 \leq t_2$ — standard temporal ordering.  If
$\Lambda_\psi = 0$, then $\psi$ must remain constant along any causal curve,
effectively decoupling it.
\end{remark}

\begin{definition}[Causal curve in complex time]
A smooth path $\gamma: [0,1] \to \mathbb{R}^4 \times \mathbb{C}_\tau$,
$\gamma(s) = (x^\mu(s), \tau(s))$, is \emph{causal} if
\begin{enumerate}
  \item $g_{\mu\nu}\,\dot{x}^\mu \dot{x}^\nu \leq 0$ (non-spacelike in the
    real metric), and
  \item $\operatorname{Re}(\dot\tau) \geq 0$ (non-decreasing real time
    component).
\end{enumerate}
\end{definition}

%----------------------------------------------------------------------
\section{Null, Timelike, and Spacelike Separation with Complex $\tau$}
%----------------------------------------------------------------------
The UBT proper-distance element generalises the Lorentzian interval.
In the complex-time sector we write
\begin{equation}
  d\mathbf{s}^2 \;=\;
  g_{\mu\nu}\,dx^\mu dx^\nu
  + \lambda_\tau\,|d\tau|^2,
\end{equation}
where $\lambda_\tau \in \mathbb{C}$ encodes the coupling of complex time to
the real geometry.  Separating real and imaginary parts:
\begin{equation}
  d\mathbf{s}^2 \;=\;
  \underbrace{g_{\mu\nu}\,dx^\mu dx^\nu
  + \operatorname{Re}(\lambda_\tau)(dt^2 - d\psi^2)}_{\text{real part}}
  + i\;\underbrace{[\cdots]}_{\text{imaginary part}}.
\end{equation}

\begin{definition}[Extended null condition]
An infinitesimal displacement $(dx^\mu, d\tau)$ is \emph{null in the
extended sense} if
\begin{equation}
  \operatorname{Re}(d\mathbf{s}^2) = 0
  \quad\text{and}\quad
  \operatorname{Im}(d\mathbf{s}^2) = 0.
\end{equation}
\end{definition}

Under Interpretation B (phase coordinate), the imaginary part vanishes
identically for real-sector observers, and the null condition reduces to
the standard Lorentzian condition $g_{\mu\nu}\,dx^\mu dx^\nu = 0$.

%----------------------------------------------------------------------
\section{Physical Consequences of Each Interpretation}
%----------------------------------------------------------------------
\begin{center}
\renewcommand{\arraystretch}{1.4}
\begin{tabular}{|l|p{4.2cm}|p{4.2cm}|p{4.2cm}|}
\hline
 & \textbf{A: Second time} & \textbf{B: Phase coord.} & \textbf{C: Auxiliary}\\
\hline
Causal dimension & $(1+1)+3 = 5$-dim & $1+3$ (standard) & $1+3$ (standard)\\
Arrow of time & Two independent arrows & Single real-$t$ arrow & Single real-$t$ arrow\\
CTC potential & High (wrapping in $(t,\psi)$) & Indirect (via $g_{\mu\nu}$ deformation) & None\\
Unitarity & Requires extra constraint & Automatic in real sector & Automatic\\
GR recovery ($\psi\to 0$) & Non-trivial & Smooth & Trivial\\
\hline
\end{tabular}
\end{center}

%----------------------------------------------------------------------
\section{Conjecture on Causal Consistency}
%----------------------------------------------------------------------
\begin{conjecture}[Causal consistency under Interpretation B]
\label{conj:causal_consistency}
Under Interpretation B (internal phase), the canonical UBT field equation
$\mathcal{E}_{\mu\nu} = \kappa\,\mathcal{T}_{\mu\nu}$ does not generate
causal paradoxes in the real sector: the real projection $g_{\mu\nu}$ of
any solution satisfies the standard Lorentzian causality conditions.
The imaginary sectors $h_{\mu\nu}$, $a_{\mu\nu}$, $b_{\mu\nu}$ modify
the causal structure only through their contribution to $T_{\mu\nu} :=
\operatorname{Re}(\mathcal{T}_{\mu\nu})$, which acts as an effective
stress-energy source inside the standard Einstein equations.
\end{conjecture}

\begin{remark}
Conjecture~\ref{conj:causal_consistency} may not hold under
Interpretation A.  Whether Interpretation A is physically admissible
within UBT remains an open question requiring analysis of the full
biquaternionic field equations; see \texttt{ctc\_conditions.tex} and
\texttt{chronology\_protection\_ubt.tex}.
\end{remark}

%----------------------------------------------------------------------
\section{Summary}
%----------------------------------------------------------------------
Three interpretations of the imaginary time component $\psi$ have been
examined.  Interpretation B (internal phase coordinate) is the most
conservative and is consistent with the canonical UBT programme: it
preserves the single real-time causal arrow and reduces smoothly to GR
in the $\psi \to 0$ limit.  Interpretations A and C represent more
extreme positions and carry correspondingly larger theoretical risks.

A causal partial order has been proposed for the complex time plane under
Interpretation B, and the notions of null, timelike, and spacelike
separation have been extended to complex $\tau$.

\medskip
\noindent\textbf{All statements in this document are speculative.  No
claim constitutes an established result of the canonical UBT framework.}

%----------------------------------------------------------------------
\bibliographystyle{plain}
\begin{thebibliography}{9}
\bibitem{Hawking1992}
  S.~W.~Hawking,
  \emph{Chronology protection conjecture},
  Phys.\ Rev.\ D \textbf{46}, 603 (1992).
\bibitem{UBT_core}
  D.~Jaroš,
  \emph{Unified Biquaternion Theory: core field equations},
  Repository document \texttt{consolidation\_project/ubt\_core\_main.tex} (2025).
\end{thebibliography}

\end{document}
